\section{Interfacet}
\label{sec:interface}

\begin{cppcode}
namespace driver::can
{
struct Frame;

class Interface
{
public:
    virtual ~Interface() noexcept = default;

    [[nodiscard]] virtual bool send(const Frame& frame)      = 0;
    [[nodiscard]] virtual bool receive(Frame& frame)         = 0;
    [[nodiscard]] virtual bool hasError() const noexcept     = 0;
    virtual void clearError() noexcept                       = 0;
};
} // namespace driver::can
\end{cppcode}

Varje del av den, en i taget:

\textbf{Virtuell destruktor.} Objekt kommer att raderas genom en \code{Interface*}. Utan
\code{virtual} anropas bara basklassens destruktor, och en \code{Spi} som håller resurser läcker
dem --- odefinierat beteende, och en av de få buggarna i den här designen som inget test fångar.
\code{= default} därför att det inte finns något att göra; \code{noexcept} därför att destruktorer
aldrig ska kasta.

\textbf{\code{= 0} på alla fyra.} Klassen är ren abstrakt: den går inte att instansiera och har
ingen implementation att ärva av misstag. Den är ett kontrakt och ingenting annat.

\textbf{\code{[[nodiscard]]} på tre av fyra.} \code{send()} returnerar \code{false} när hårdvaran
är upptagen. En anropare som ignorerar det tror att framen gick iväg --- den vanligaste
drivrutinsbuggen som finns, här förvandlad till en kompilatorvarning.

\textbf{\code{const} på \code{hasError()}.} Att fråga om ett fel ändrar inte drivrutinen, och
metoden går därmed att anropa på en \code{const Interface&}.

\textbf{\code{noexcept} på de två felmetoderna.} De anropas från felhanteringsvägar, och en
felhanteringsväg som själv kan kasta är svår att resonera om. De gör dessutom ingenting som
\emph{kan} misslyckas.

\subsection{Vad som inte finns här}

Ingen \code{init()}. Ingen \code{setBaudRate()}. Ingen \code{readRegister()}.

\textbf{Testet är: kan både \code{Stub} och \code{Spi} implementera metoden meningsfullt?}
\code{readRegister()} kan de inte --- stubben har inga register. En sådan metod hade tvingat
stubben att ljuga, och tvingat applikationen att veta att register finns.

\section{Stubben}
\label{sec:stub}

\code{Stub} är \code{final}, ärver \code{Interface}, och simulerar CAN i minnet:

\begin{booktable}{@{}lL@{}}
\thead{Metod} & \thead{Kontrakt} \\ \midrule
Konstruktorn & \code{myLastFrame} tom, \code{myInitialized = true}, \code{myHasData = false}. \\
\code{send()} & \code{false} om oinitialiserad. Annars: lagra framen, sätt \code{myHasData},
  returnera \code{true}. \\
\code{receive()} & \code{false} om oinitialiserad eller utan data. Annars: kopiera ut framen,
  nollställ \code{myHasData}, returnera \code{true}. \\
\code{hasError()} & Alltid \code{false}. \\
\code{clearError()} & Gör ingenting. \\
\code{inject()} & Lagrar en frame för senare \code{receive()}. \\
\end{booktable}

Kopierings- och flyttoperationerna tas bort i den publika delen. En \code{Stub} representerar en
(simulerad) hårdvaruresurs, och två kopior av samma resurs är inte två resurser.

\subsection{Varför \texorpdfstring{\code{inject()}}{inject()} inte ligger i interfacet}

Den finns bara för testernas skull: den simulerar att en frame kommit in på bussen. En \code{Spi}
kan inte implementera den --- den kan inte trolla fram trafik som inte finns.

Konsekvensen är mönstret som återkommer genom hela boken:

\begin{cppcode}
driver::can::Stub stub{};            // konkret typ: har inject()
app::EchoNode node{stub, 0x200U};    // tar en Interface&: har den inte
stub.inject(frame);                  // testet talar med stubben direkt
node.run();                          // komponenten vet ingenting om det
\end{cppcode}

Testet håller i den \textbf{konkreta} typen, komponenten ser bara \textbf{interfacet}.

\subsection{Vad stubben inte simulerar}

Inte fel. Inte upptagen hårdvara. Ingen kö --- en andra \code{send()} skriver över den första.
Ingen arbitrering, ingen CRC.

Det är avsiktligt. Stubbens jobb är att låta \code{EchoNode} köras, inte att vara en
CAN-simulator.

\begin{aside}
\note Att veta var gränsen för en testdubblares realism går är en färdighet i sig. \textbf{En
stubb som växer tills den är en andra implementation av hårdvaran är en stubb som själv behöver
testas.} Den som vill testa hur drivrutinen beter sig när hårdvaran är upptagen behöver en annan
dubblare, en söm längre ned --- och den kommer i \kapref{ch:testning}.
\end{aside}

\subsection{Stubb, inte mock}

En \textbf{stubb} svarar på anrop med förprogrammerade svar; testet kontrollerar vad den anropande
koden gjorde med svaren. En \textbf{mock} kontrollerar dessutom att den blev anropad rätt --- rätt
metod, rätt argument, rätt antal gånger.

\code{Stub} är en stubb. \code{ScriptedTransport} i \kapref{ch:testning} ligger närmare en mock.

\subsection{Sömmen har redan lönat sig}

Argumentet för ett interface brukar vara hypotetiskt. Här är det inte det: en tidigare version av
konstruktionen nådde registren via \textbf{minnesmappad I/O} på \code{0xFF200000}, vilket
förutsätter en adressbuss ut mot FPGA-väven som mikrokontrollern inte har.

Allt ovanför \code{Interface} överlevde bytet till SPI \textbf{oförändrat}. Den gamla varianten
hade \code{Fpga} och \code{volatile}-pekare där vi har \code{Spi} och \code{ByteTransport}. Ingen
av dem är fel; det är två implementationer under samma söm.
